\scanpage{51}
% archive foliation: 50
(ii) Ce raisonnement vaut aussi sans supposer $T/S$ affine, si $g$ est \uline{étale}, \deleted{\unsure{car}} Ainsi, $f_*(g)$ sera simple. Pour prouver qu'il est étale, il suffit de prouver que ses \uline{fibres sont} \deleted{\unsure{de}} \inserted{finies} \deleted{dim $0$ (\unsure{uniformément})}. \inserted{(\unsure{uniformes})} On est ramené au cas où $S=\operatorname{Spec}(k)$, $k$ un corps, et à prouver que si $k'$ est une clôture algébrique de $k$, alors il n'y a qu'un nb fini de $k'$-pts de $f_*(X)$ au dessus d'un $k'$-point de $f_*(Y)$, i.e. qu'il n'y a qu'un nb fini de sections de $X'/T'$ qui relèvent une section donnée de $Y'/T'$.
\margin{\unsure{donner un} \uline{\unsure{contre-ex.}} \unsure{non fini}}
\diagram{\begin{tikzcd}&X \arrow[r]&Y\\ k&T\end{tikzcd}}
\diagram{\begin{tikzcd}&f_*(X)\arrow[d]\\\operatorname{Spec}(k')\arrow[r]&f_*(Y)\end{tikzcd}}

[$\unsure{\text{où}}$ $'$ désigne ext. des scalaires de $k$ à $k'$]
Or $T'$ sera noethérien [de type fini sur $k'$] donc n'a qu'un nb fini de comp. connexes. On peut supposer $T'$ connexe. Mais le relèvement est connu dès qu'on connaît l'image d'un pt. On \unsure{trouve|aura} donc

\uline{Corollaire} 1 Soit $\nu$ le maximum des nbs \deleted{\unsure{resp.}} de composantes connexes absolues des $T_s$ ($s\in S$), et soit de $=$ $n$ le maximum analogue pour $X/Y$. Alors le maximum
